% Fixture "plain": standard article, one column, default margins.
% Exercises ligatures (fi, ff, ffl), umlauts, possessives, line-end hyphenation, inline math,
% a phrase that occurs twice, and a second page.
\documentclass[11pt]{article}
\usepackage[T1]{fontenc}
\usepackage[utf8]{inputenc}
\usepackage{amsmath,amsthm}
\newtheorem{theorem}{Theorem}
\pdfinfoomitdate=1 \pdftrailerid{} \pdfsuppressptexinfo=-1
\title{On Continuous Maps of the Interval}
\author{Anna Müller}
\date{}
\begin{document}
\maketitle

\section{Introduction}
The map $f$ is continuous on the closed interval $[0,1]$ and differentiable on its interior.
We study the efficient computation of fixed points; the affine case is trivial, but the general
one is not. In Müller's earlier work the fluffy notion of a ``nice'' function was used without
a definition, which we now remedy. The map $f$ is continuous on the closed interval only when
$f(0)$ is defined, a subtlety that is easy to overlook in a first reading of the argument.

\begin{theorem}
Every continuous map $f\colon [0,1]\to[0,1]$ has a fixed point.
\end{theorem}
\begin{proof}
Let $g(x) = f(x) - x$. Then $g(0) \ge 0$ and $g(1) \le 0$, so by the intermediate value theorem
there is $x^\ast$ with $g(x^\ast) = 0$.
\end{proof}

A considerably longer sentence follows here so that the typesetter is forced to hyphenate some
extraordinarily long words such as incomprehensibilities, characteristically, and
internationalization near the right margin of the paragraph.
\[
  \int_0^1 f(x)\,dx = \sum_{n\ge 1} a_n, \qquad a_n = \frac{1}{n^2}.
\]

\newpage
\section{Second page}
This page exists to test page ranges and anchors that are not on the first page. The map $f$ is
mentioned once more here.
\end{document}
